\documentclass[10pt,twocolumn]{article}
\usepackage[T1]{fontenc}
\usepackage[utf8]{inputenc}
\usepackage{lmodern}
\usepackage[margin=1.8cm,top=2.4cm,bottom=2cm,columnsep=0.6cm]{geometry}
\usepackage{fancyhdr}
\usepackage{titlesec}
\usepackage{booktabs}
\usepackage{amsmath,amssymb,amsfonts}
\usepackage{microtype}
\usepackage{xcolor}
\usepackage{mdframed}
\usepackage{parskip}
\usepackage{enumitem}
\usepackage{hyperref}

\definecolor{rulered}{RGB}{180,30,30}
\definecolor{boxgray}{RGB}{245,245,245}
\definecolor{keywordgray}{RGB}{100,100,100}

\pagestyle{fancy}
\fancyhf{}
\fancyhead[L]{\textsc{\small Claude Mag}}
\fancyhead[C]{\small\textit{Machine Learning Research Digest}}
\fancyhead[R]{\small 2026-10-26}
\fancyfoot[C]{\small\thepage}
\renewcommand{\headrulewidth}{0.8pt}

\titleformat{\section}{\normalsize\bfseries\color{rulered}}{}{0pt}{}%
  [\vspace{-4pt}\color{rulered}\rule{\columnwidth}{0.4pt}]
\titlespacing{\section}{0pt}{8pt}{4pt}

\hypersetup{colorlinks=true,urlcolor=rulered,linkcolor=black}

\begin{document}
\setlength{\parindent}{0pt}
\setlength{\parskip}{4pt}

{\small\color{keywordgray}\textbf{Keywords:}
  text-to-image generation \textbullet{}
  compositional reasoning \textbullet{}
  iterative refinement \textbullet{}
  GRPO}\par\vspace{4pt}

{\LARGE\bfseries\raggedright
  FiRe: Fine-grained Multimodal Reasoning for Enhanced Image
  Generation}\par\vspace{4pt}

{\large\itshape\raggedright
  Requirement-Level Check-Refine Cycles and Per-Step RL Rewards
  Fix the Fix-One-Break-Another Problem in Compositional T2I}\par\vspace{6pt}

{\small\textbf{By} Yongjin Kim, Yoonjin Oh, Yerin Kim, Hyomin Kim,
  Jeeyoung Yun, Yujung Heo, Minjun Kim, Sungwoong Kim (KU-AGI)}\par\vspace{2pt}

{\small\color{keywordgray}arXiv:2604.13491 \textbullet{}
  NeurIPS 2026}
\par\vspace{2pt}\noindent{\color{rulered}\rule{\columnwidth}{0.6pt}}\vspace{6pt}

Text-to-image models fail on compositional prompts because whole-image
feedback cannot isolate which individual visual requirement is
violated---so fixing one element inadvertently breaks another. FiRe
decomposes prompts into atomic visual requirements, checks each one
independently using a multimodal LLM, and applies \emph{local}
refinements to only the failing regions. FiRe-GRPO, the accompanying
reinforcement learning variant, assigns a per-step reward to each
check-refine action, enabling the model to learn when and how to refine
rather than merely whether the final image is correct. On compositional
benchmarks including T2I-CompBench and GenAI-Bench, FiRe achieves
consistent improvement over T2I baselines, with the largest gains on
prompts specifying three or more simultaneous visual constraints.

\begin{mdframed}[backgroundcolor=boxgray,linecolor=rulered,linewidth=1pt,
    innerleftmargin=6pt,innerrightmargin=6pt,
    innertopmargin=6pt,innerbottommargin=6pt]
\textbf{\small AT A GLANCE}\par\vspace{4pt}
\begin{description}[leftmargin=0pt,labelsep=4pt,itemsep=2pt,parsep=0pt,
    font=\small\bfseries]
  \item[Problem:] Compositional T2I prompts impose multiple simultaneous
    visual constraints; whole-image feedback cannot localize failures or
    guide targeted corrections.
  \item[Why it matters:] State-of-the-art T2I models still fail on multi-constraint
    prompts despite improved language understanding; the bottleneck is the
    feedback signal, not the generator.
  \item[Method:] Decompose the prompt into visual requirements; check each
    with an MLLM; refine locally; train with FiRe-GRPO assigning step-level
    rewards to each check-refine action.
  \item[Result:] Consistent improvement on compositional benchmarks,
    especially on prompts with $\geq 3$ simultaneous constraints.
  \item[Evidence:] T2I-CompBench, GenAI-Bench; ablations isolate local vs.\
    global refinement and step-level vs.\ terminal-only reward.
\end{description}
\end{mdframed}

\section{The Big Picture}

Language-guided image generation has improved dramatically, yet compositional
prompts---requests that simultaneously specify object counts, attributes,
spatial relationships, and interactions---remain a known failure mode. A
prompt such as ``a red cup to the left of a blue vase on a wooden table''
imposes three simultaneous visual constraints; any failure in the generated
image requires identifying \emph{which} constraint is violated before attempting
a correction.

Existing reasoning-augmented approaches attempt to fix this by elaborating the
prompt before generation or by scoring the whole image after generation. Neither
approach can localize the failure to a specific requirement or avoid disrupting
already-satisfied constraints during refinement.

FiRe reframes image generation as a sequential decision process where each step
either checks one requirement or applies one local refinement, and learns via
reinforcement learning to execute this process efficiently.

\section{Why Existing Approaches Fall Short}

\textbf{Prompt augmentation (chain-of-thought elaboration):}
Improves average image quality but generates the same elaborated prompt for
every seed; it cannot recover from specific failures in a generated image.

\textbf{Whole-image feedback (DALL-E Reflect, ImageReward):}
Produces a scalar or sentence-level critique. This is too coarse to guide
targeted corrections and does not indicate which region or property to change.

\textbf{Terminal-only RL for T2I:}
Assigning reward only to the final image creates a sparse signal that is
difficult to optimize and does not credit or penalize intermediate refinement
decisions.

\section{Core Mechanism}

\textbf{Step 1 --- Requirement decomposition.}
An MLLM parses the prompt into a list of atomic visual requirements
$\mathcal{R} = \{r_1, \ldots, r_n\}$, where each $r_i$ is a checkable
property (e.g., ``red cup present'', ``cup to the left of the vase'').

\textbf{Step 2 --- Initial generation.}
A base T2I model generates $I_0$ from the prompt.

\textbf{Step 3 --- Fine-grained check.}
For each unsatisfied requirement $r_i$, the MLLM answers ``Is $r_i$
satisfied in $I_{t-1}$?'' with a binary verdict and a natural-language
diagnosis specifying the region and nature of the failure.

\textbf{Step 4 --- Local refinement.}
The least-satisfied requirement $i^* = \arg\min_i \mathcal{V}_i(I_{t-1})$
triggers an inpainting or latent edit confined to the region described in the
diagnosis, yielding $I_t$. Unaffected regions are left unchanged.

\textbf{Step 5 --- Iteration.}
Steps 3--4 repeat until all requirements are satisfied or the step budget
$T$ is exhausted.

\section{Technical Formulation}

Let $\mathcal{V}_i(I) \in [0,1]$ denote the MLLM's estimated probability
that requirement $r_i$ is satisfied in image $I$. The step-level reward at
refinement step $t$ (targeting requirement $i^*$) is:
\[
\rho_t = \mathcal{V}_{i^*}(I_t) - \mathcal{V}_{i^*}(I_{t-1})
\]
The FiRe-GRPO policy gradient objective is:
\[
\mathcal{J}_{\text{GRPO}}(\theta)
  = \mathbb{E}\!\left[\sum_{t=1}^T \rho_t \cdot \log \pi_\theta(\delta_t \mid I_{t-1}, r_{i^*})\right]
  - \beta \cdot D_{\text{KL}}(\pi_\theta \| \pi_{\text{ref}})
\]
where $\delta_t$ is the refinement action (diagnosis + edit), $\pi_\theta$ is
the MLLM policy being trained, and $\pi_{\text{ref}}$ is a frozen reference
policy. The KL term prevents the policy from diverging far from a well-behaved
prior and controls mode collapse.

The combined requirement satisfaction score at the end of generation is:
\[
\text{Score}(I_T) = \frac{1}{n} \sum_{i=1}^n \mathcal{V}_i(I_T)
\]
This is the final metric reported on benchmarks.

\section{Learning or Inference Procedure}

\textbf{Training:}
\begin{enumerate}[itemsep=2pt,parsep=0pt]
  \item Sample a compositional prompt; decompose into requirements.
  \item Generate $I_0$ with the base T2I model.
  \item Run the check-refine cycle for $T$ steps; collect $\{\rho_t\}$.
  \item Update the MLLM backbone with FiRe-GRPO.
  \item The 300-step checkpoint (FiRe-300Step) is used for all evaluations.
\end{enumerate}

\textbf{Inference:}
\begin{enumerate}[itemsep=2pt,parsep=0pt]
  \item Parse the prompt; generate $I_0$.
  \item Iterate: check all requirements, identify the most violated,
    apply a local refinement.
  \item Return $I_T$ after the budget is exhausted or all requirements
    are satisfied.
\end{enumerate}

\section{What the Guarantee Says}

The step-level reward $\rho_t$ provides a per-requirement improvement signal.
Under mild assumptions on the MLLM's calibration (i.e., $\mathcal{V}_i$
monotonically correlates with human satisfaction judgments), the GRPO
objective with the KL regularizer converges to a policy that maximizes
expected total requirement satisfaction while remaining close to the
reference policy. This ensures that FiRe does not sacrifice the reference
model's general image quality in pursuit of compositional accuracy.

\section{Experimental Findings}

\begin{itemize}[itemsep=2pt,parsep=0pt]
  \item \textbf{T2I-CompBench:} FiRe improves attribute binding and spatial
    relationship accuracy over all baselines; largest gains on prompts with
    $\geq 3$ simultaneous constraints.
  \item \textbf{GenAI-Bench:} Consistent improvement across compositional
    categories (color, count, spatial, interaction).
  \item \textbf{Fix-one-break-another:} Global refinement baselines show
    improvement in the targeted attribute but regression in other attributes;
    FiRe's local edits preserve already-satisfied requirements.
  \item \textbf{FiRe-GRPO vs.\ supervised fine-tuning:} Step-level RL
    converges faster and produces better refinement policies than training
    with supervised labels for each step.
\end{itemize}

\section{Ablations and Interpretation}

\textbf{Global vs.\ local refinement:}
Applying refinement to the whole image rather than the failing region degrades
already-satisfied requirements by an average of 12 percentage points on
T2I-CompBench; local refinement avoids this regression entirely.

\textbf{Step-level vs.\ terminal-only reward:}
Terminal-only reward (only $\mathcal{V}_{i^*}(I_T)$ at the final step)
trains 40\% slower and achieves 8 percentage points lower compositional
accuracy; step-level rewards provide a denser credit assignment signal.

\textbf{Number of refinement steps $T$:}
Performance improves with $T$ up to 3--4 steps and plateaus beyond; additional
steps add compute without further gains on typical 3-constraint prompts.

\textbf{Without requirement decomposition:}
Using a single holistic verdict instead of per-requirement checks degrades
performance to below the whole-image feedback baseline, confirming that
the decomposition is the key structural contribution.

\section*{Reference}

Yongjin Kim, Yoonjin Oh, Yerin Kim, Hyomin Kim, Jeeyoung Yun, Yujung Heo,
Minjun Kim, Sungwoong Kim.
\textbf{FiRe: Fine-grained Multimodal Reasoning for Enhanced Image Generation}.
arXiv:2604.13491. NeurIPS 2026.
\url{https://arxiv.org/abs/2604.13491}

\end{document}
